%! Rational Functions and Vertical Asymptotes — Unit 3, Lesson 3.3 — Exit Ticket (Answer Key)
\documentclass[10pt]{article}
\usepackage{precalculus-article}
\usepackage{precalculus-key}
\newcommand{\TallMath}[1]{$\displaystyle #1\rule[-1.4em]{0pt}{3.2em}$}

\begin{document}

\pageheader{Unit 3, Lesson 3.3}{Exit Ticket --- Answer Key}
\namedateperiod

\vspace{0.14in}

\begin{enumerate}[itemsep=16pt]

\item For $r(x) = \dfrac{x - 1}{(x + 2)(x - 4)}$:

\begin{enumerate}[label=(\alph*), itemsep=8pt]
  \item Find all values of $x$ where the denominator equals zero.

        \ans{$(x+2)(x-4) = 0 \Rightarrow x = -2$ and $x = 4$}

  \item For each value from part (a), evaluate the numerator to determine whether it
        gives a vertical asymptote.

        At $x = -2$: numerator $=$ \ans{$-2-1 = -3 \ne 0$}. \quad
        This gives: \ans{vertical asymptote at $x = -2$}

        At $x = 4$: numerator $=$ \ans{$4-1 = 3 \ne 0$}. \quad
        This gives: \ans{vertical asymptote at $x = 4$}

  \item Write the equation(s) of all vertical asymptotes of $r(x)$.

        \ans{$x = -2$ and $x = 4$}
\end{enumerate}

\item Write the one-sided limits.

\[
  \lim_{x \to 3^+} r(x) = \ans{$+\infty$}
  \qquad\qquad
  \lim_{x \to 3^-} r(x) = \ans{$-\infty$}
\]

\end{enumerate}

\end{document}
